% !TeX program = lualatex
% Copyright (C) 2026 IEEE UFG Computer Society contributors.
% Distributed under the LaTeX Project Public License 1.3c or later.
% See LICENSE and development/MANIFEST.md.
%
% MODELO DE RELATÓRIO DE ATIVIDADE
% Use após uma reunião, evento, ação de extensão ou atividade de projeto.
% Registre a execução e as evidências; decisões formais ficam na ata.
% Substitua todos os campos entre colchetes e remova os não aplicáveis.
% Os sistemas do IEEE e da UFG mantêm seus próprios registros oficiais.
\documentclass[
  paper-size=a4,
  font-size=11pt,
  two-sided=false,
  title-page=false,
  section-numbering=plain,
  inner-margin=22mm,
  outer-margin=22mm,
  top-margin=24mm,
  bottom-margin=28mm,
  color-model=rgb
]{ieeeufgcs}
\usepackage[brazilian]{babel}
\usepackage{indentfirst}
\usepackage[hidelinks,hypertexnames=false]{hyperref}
\raggedbottom

% Edite estes valores uma vez; eles aparecem no título e na identificação.
\newcommand{\RelatorioAtividade}{[nome da atividade]}
\newcommand{\RelatorioEmissao}{[data de emissão]}

\IEEECSSetup{
  document-type={Relatório de atividade},
  title={\RelatorioAtividade},
  subtitle={Capítulo Estudantil IEEE Computer Society UFG},
  running-title={Relatório de atividade},
  institution={Capítulo Estudantil da IEEE Computer Society na UFG},
  author={[Nome e função de quem redigiu]},
  date={\RelatorioEmissao},
  location={Goiânia},
  identifier={Minuta para revisão}
}

\begin{document}
\maketitle

\begin{IEEECSBrandedTable}{
  >{\raggedright\arraybackslash}m{0.30\linewidth}
  >{\raggedright\arraybackslash}X
}
\IEEECSHeaderCell{Identificação} & \IEEECSHeaderCell{Registro} \\
Atividade & \RelatorioAtividade. \\
Tipo & [Reunião, evento, ação de extensão ou atividade de projeto]. \\
Data e horário & [Dia, início e término]. \\
Local ou modalidade & [Unidade, sala ou plataforma]. \\
Equipe responsável & [Área, coordenação e demais responsáveis]. \\
Público & [Público pretendido e participantes efetivos]. \\
Ata relacionada & [Número ou link, quando houver deliberações]. \\
Situação & [Minuta, aprovado ou arquivado]. \\
\end{IEEECSBrandedTable}

\section*{Objetivo e preparação}

[Descreva o propósito da atividade, a necessidade atendida e o resultado
esperado. Informe a divulgação, as inscrições, as autorizações ou os
materiais preparados que sejam relevantes para avaliar a execução.]

\section*{Execução}

[Registre o que ocorreu, em ordem lógica, com datas e responsáveis quando
necessário. Diferencie o que estava previsto do que foi realizado. Para
reuniões, cite a ata que contém as deliberações.]

\section*{Resultados e avaliação}

\begin{IEEECSBrandedTable}{
  >{\raggedright\arraybackslash}m{0.32\linewidth}
  >{\raggedright\arraybackslash}X
}
\IEEECSHeaderCell{Indicador ou entrega} &
\IEEECSHeaderCell{Resultado e fonte} \\
\emph{[Presença ou alcance]} &
\emph{[Quantidade apurada e fonte da contagem]}. \\
\emph{[Entrega técnica ou formativa]} &
\emph{[Resultado verificável e localização]}. \\
\emph{[Retorno dos participantes]} &
\emph{[Método de coleta e síntese, se houver]}. \\
\end{IEEECSBrandedTable}

[Registre o que funcionou, as limitações, os incidentes e as lições
aprendidas. Evite atribuir resultados que não possam ser verificados.]

\section*{Recursos e evidências}

\textbf{Recursos usados:} [Espaço, equipamentos, materiais e apoio recebido].

\textbf{Despesas e receitas:} [Valores realizados ou ``sem movimentação'';
indique onde estão os comprovantes e a prestação de contas].

\textbf{Evidências:} [Links ou caminhos de lista de presença, materiais,
fotografias autorizadas, repositórios, formulários e resultados].

% Guarde dados pessoais e comprovantes em local de acesso apropriado.
% Descreva aqui a evidência, sem copiar dados sensíveis para o PDF.

\section*{Pendências e encaminhamentos}

\begin{IEEECSBrandedTable}{
  >{\raggedright\arraybackslash}X
  >{\raggedright\arraybackslash}m{0.24\linewidth}
  >{\raggedright\arraybackslash}m{0.17\linewidth}
}
\IEEECSHeaderCell{Ação} &
\IEEECSHeaderCell{Responsável} &
\IEEECSHeaderCell{Prazo} \\
\emph{[Próxima ação ou ``nenhuma'']} &
\emph{[Nome ou função]} & \emph{[Data]} \\
\end{IEEECSBrandedTable}

\section*{Revisão e registros oficiais}

\begin{IEEECSBrandedTable}{
  >{\raggedright\arraybackslash}m{0.30\linewidth}
  >{\raggedright\arraybackslash}X
}
\IEEECSHeaderCell{Etapa} & \IEEECSHeaderCell{Responsável e registro} \\
Redação & [Nome, função e data]. \\
Revisão ou aprovação & [Nome, função, data e decisão]. \\
IEEE vTools, se aplicável & [Identificador e situação do envio]. \\
UFG SIGAA, se aplicável & [Identificador e situação do cadastro]. \\
Arquivo final & [Local da versão aprovada e dos anexos]. \\
\end{IEEECSBrandedTable}

% Após o preenchimento, retire linhas de sistemas que não se aplicam.
% Confirme os requisitos atuais nos sistemas oficiais antes do envio.

\end{document}
